\documentclass[a4paper,12pt]{article}
\usepackage{amsmath, amssymb}
\usepackage{xcolor}
\usepackage{ulem}        % underline
\usepackage{tcolorbox}   % colored boxes (if needed)
\usepackage{pifont}
\usepackage{xcolor}
\usepackage{tikz}
\usepackage{centernot}
\usepackage{booktabs}
\usepackage{esint} % Extended integrals
\renewcommand{\notin}{\mathrel{\not\in}}
\usetikzlibrary{trees} % for tree structures

\newcommand{\Eemp}{\hat{\mathbb{E}}}
\newcommand{\Varhat}{\widehat{\operatorname{Var}}}
\newcommand{\shat}{\hat{\sigma}}
\newcommand{\Covhat}{\widehat{\operatorname{Cov}}}
\newcommand{\chat}{\hat{c}}
\newcommand{\corrcoef}{\hat{r}}


\begin{document}

\begin{center}
\begin{center}
    \huge{\textbf{Mathématiques : CM}}
\end{center}
\end{center}
\newpage

\tableofcontents 


\newpage



\section{Sequences of functions}

Let $I$ be an interval of $\mathbb{R}$ ($f_n$) a sequence of functions on $I$. \newline

Sequence whose general term $f_n$ is a function $f_n \, : I \rightarrow \mathbb{R}$. \newline

Modes of convergence of ($f_n$) to $f$ where $f \, : \rightarrow \mathbb{R}$. \newline

Pointwise convergence :

($f_n$) converges pointwise to a function $f$ if  $\forall \, x \, \in I$, $f_n(x) \underset{n \rightarrow +\infty}{\rightarrow} f(x)$

Example : on $I = [0, 1], f_n(x) = x^n$

\[
\text{Let } x\in [0,1], \begin{cases}
    \text{if } x < 1, x^n \rightarrow 0 \\
    \text{if } x = 1, x^n \rightarrow 1
\end{cases}
\]

Thus $f_n(x) \underset{n \rightarrow +\infty}{\rightarrow} f(x)$ with :

\[
f(x) = \begin{cases}
    0 \text{ if } x < 1 \\
    1 \text{ if } x = 1
\end{cases}
\]

\[(f_n) \overset{\text{p.w}}{\rightarrow} f \text{ (i.e converges pointwise)} \]

\section{Uniform convergence}

Supremum of a bounded set $S$. \newline
Let $S$ be a set, upperbounded. \newline

$\exists M \in \mathbb{R}$, $\forall x \in S$, $x \leq M$ \newline

Supremum of $S$ : the smallest possible $M$ denoted $Sup S$. \newline

Example : $S = [0,3], Sup S = 3$ \newline
$S_2 = [0,3[, Sup S = 3$ \newline

Remark : For $[0,3]$, $Sup [0,3] = 3 = max[0,3]$ \newline

For $[0,3[$ = no maximum value ($3 \notin [0,3[$) but $Sup [0,3] = 3$

\section{Definition : norm $||.||_\infty$}

Let $E$ be the vector space of bounded function on $I$. \newline

Let $f \in E$,

TODO 

\subsection{Properties}

\[
\forall (f,g) \in E^2, \forall \alpha \in \mathbb{R} \\
\]

\[
||f||_\infty = 0 \implies f = 0_E 
\]

\[
||\alpha f||_\infty = ||\alpha|| \times ||f||_\infty
\]

\[
||f + g||_\infty \leq ||f||_\infty + |||g|_\infty
\] \newline

$(f_n)$ converges uniformly to f if $\forall n \geq n_0, f_n - f$ is bounded $\implies ||f_n - f||_\infty \rightarrow 0$

Example : on $[0,1], f_n(x) = x^n$. \newline

We have seen that : $\forall x \in [0,1], f_n(x) \rightarrow f(x) = \begin{cases}
    0 \text{ if } x < 1 \\
    1 \text{ if } x = 1
\end{cases}$

\[(f_n) \overset{\text{p.w}}{\rightarrow} f \text{ (i.e converges pointwise)} \] 

Does $(f_n)$ converges uniformly to $f$ ? \newline

Let $n \in \mathbb{N}$. Let us study $f_n - f$ : \newline

\[
\forall x \in [0, 1], f_n(x) - f(x) = \begin{cases}
    x^n \text{ if } x < 1\\
    0 \text{ if } x = 1
\end{cases}
\]

\[
f_n(1) - f(1) = 1^n - 1 = 0
\]

Thus $(f_n)$ converges uniformly. \newline

Example 2 : on $\mathbb{R}, f_n(x) = \frac{sin(nx)}{n}$ \newline

Pointwise : let $x\in \mathbb{R}$ \newline

\[
\forall n \in \mathbb{N}, -\frac{1}{n} \leq \frac{sin(nx)}{n} \leq \frac{1}{n}
\]

\[-\frac{1}{n} \rightarrow 0 \land \frac{1}{n} \rightarrow 0\]

Thus, $f_n(x) \rightarrow 0$ \newline

$f_n(x) \overset{p.w}{\rightarrow} f$ (zero function) \newline

$\forall x \in \mathbb{R}, f(x) = 0$ \newline

Does ($f_n$) converges to $f$ ? \newline

Let $n \in \mathbb{N}^*$, study $f_n - f$ \newline

$\forall x \in \mathbb{R}, |f_n(x) - f(x)| = |\frac{sin(nx)}{n}|$ \newline

$Sup{|f_n(x) - f(x)|, x \in \mathbb{R}} = \frac{1}{n}$

\[
\frac{1}{n} \rightarrow 0 \implies ||f_n - f||_\infty \rightarrow 0 \implies (f) \overset{\text{CVG}}{\rightarrow} f
\] \newline

Thus we have a \textbf{Uniform convergence}.

\section{L2 convergence (on a closed interval [a, b]).}

\subsection{Definition - piecewise continuity}

A function $f = [a, b] \rightarrow \mathbb{R} $ is "piecewise" continuous on $[a, b]$ if it is continous at all $x \in [a, b]$, except maybe at a finite set of values $\{d_1, d_2, ..., d_n\}$ \newline

At each discontinuous point d, there is a left limit and a right limit. \newline

\[
f(d^-) = \underset{\underset{x < d}{x \rightarrow d}}{\text{lim}} \; f(x)
\]
\[
f(d^+) = \underset{\underset{x > d}{x \rightarrow d}}{\text{lim}} \; f(d)
\]

Remark : we can define $\int_{0}^{2} f(x)\,dx = \int_{0}^{1} f(x)\,dx + \int_{1}^{2} f(x)\,dx$

\subsection{Norm $||.||_2$ on $[a, b]$}

Let $E$ be the vector space of functions, $f$ piecewise continuous on $[a,b]$

\[
||f||_2 = \sqrt{\int_{a}^{b}f^2(x)\,dx}
\]

Remark :  $||.||_2$  comes from the scalar product $<f, g> = \int_{a}^{b}f(x)g(x)\,dx$ \newline

\[ 
||f||_2 = \sqrt{<f, f>}
\]

\subsection{Properties}

$\forall (f,g) \in E^2, \forall \alpha \in \mathbb{R}$ : \newline

\begin{itemize}
    \item $||f||_2 = 0 \implies f$ is the zero function (except maybe at a finite set of points)
    \item $||\alpha f||_2 = |\alpha| * ||f||_2$
    \item $||f + g||_2 \leq ||f||_2 + ||g||_2$
\end{itemize}

\subsection{Definition}

$(f_n)$ converges L2 to $f$ if : $||f_n - f||_2 \underset{n \rightarrow +\infty}{\rightarrow} 0$ \newline

Example 1 : \newline

On $[0, 1], f_n(x) = x^n$ : \newline

\[(f_n) \overset{\text{p.w}}{\rightarrow} f \text{ where }  f(x) = \begin{cases}
    0 \text{ if } x < 1 \\
    1 \text{ if } x = 1
\end{cases}\] 

\[
(f_n) \text{ does not converge to } f
\]

\[\text{Does } (f_n) \text{ converge L2 to } f \]

Let $n \in \mathbb{N}$, 

\[
f_n(x) - f(x) = \begin{cases}
    x^n \text{ if } x < 1  \\
    0 \text{ if } x = 1
\end{cases}
\]

\section{Series of functions}

\subsection{Definition}

Let $(f_n)$ be a sequence of functions. \newline

The series $\sum{f_n}$ is the sequence $(S_n)$ :

\[
\forall n \in \mathbb{N}, S_n = f_0 + f_1 + ... + f_n = \sum_{k = 0}^{n}{f_k}
\]

Example : Power series $f_n : x \mapsto a_nx^n$, $(a_n \in \mathbb{R})$ \newline

\subsection{Definition}

\begin{itemize}
    \item $\sum{f_n}$ converges point wise if $(S_n)$ converges point wise
    \item $\sum{f_n}$ converges uniformly if $(S_n)$ converges uniformly
    \item $\sum{f_n}$ converges L2 if $(S_n)$ converges L2
\end{itemize}

Necessary condition : 
\begin{itemize}
    \item $\sum f_n$ converges uniformly $\implies (f_n) \rightarrow f_n(x) = 0$ (converges uniformly)
    \item $f_n$ does not converge to the zero function $\implies \sum{f_n}$ does not converge
    \item $\sum f_n$ converges L2 $\implies (f_n) \rightarrow f_n(x) = 0$ (converges L2)
\end{itemize}

Sufficient condition :
\begin{itemize}
    \item $\sum{||f_n||_\infty}$ converge $\implies \sum{f_n}$ converge
    \item $\sum{||f_n||_2}$ converge $\implies \sum{f_n}$ converge L2
\end{itemize}

\section{Fourier Series}

Let $f$ be piecewise continuous on $\mathbb{R}$ of period $T$.
\[
\forall x \in \mathbb{R}, f(x + T) = f(x)
\]

\subsection{Case $T = 2\pi$}

\subsubsection{Definition : Fourier series of f}

\[
FS(f) = \frac{a_0}{2} + \sum_{n\geq 1}{(a_n \, cos(nx) + b_n \,sin(nx))}
\]
\[
\forall n \in \mathbb{N}, a_n = \frac{1}{\pi} \int_{-\pi}^{\pi}{f(x) \, cos(nx) \, dx}
\]

\subsubsection{Explanation}

Let $\mathbb{E}$ the $\mathbb{R}$ vector space of piecewise continuous, functions, period. 

\[
<f, g> = \frac{1}{\pi}\int_{-\pi}^{\pi}{f(x)\,g(x)\,dx}
\]

Note that :

\[
||f||^2 = <f, f> = \frac{1}{\pi}\int_{-\pi}^{\pi}{f^2(x)\,dx}
\]

($||f||$ small means $f^2$ and $|f|$ is small on average) \newline

Then : 
\[
\begin{cases}
    a_n = <f, cos(nx)> \\
    b_n = <f, sin(nx)> \\
    a_0 = <f, 1> \\
\end{cases}
\]

Let :
\[
\mathcal{F}_n = (\frac{1}{\sqrt{2}}, cos(x), sin(x), cos(2x), sin(2x), ..., cos(nx), sin(nx))
\]

$\mathcal{F}_n$ is an orthonormal family. \newline

The partial sum $S_n$ of the series is the orthogonal projection of $f$ on $Span(\mathcal{F}_n)$ defined such as:

\[
S_n = \frac{a_0}{2} + (a_1\,cos(x) + b_1\,sin(x)) + ... + (a_n\,cos(nx) + b_n\,sin(nx))
\]

$S_n$ is the element of $Span(\mathcal{F}_n)$ which minimizes $||f - S_n||^2$ \newline

$S_n$ is the function of the form :

\[
\frac{a_0}{2} + (a_1\,cos(x) + b_1\,sin(x)) + ... + (a_n\,cos(nx) + b_n\,sin(nx))
\]

$S_n$ is the element of $Span(\mathcal{F}_n)$ which minimizes $||f - S_n||^2$. \newline

$S_n$ is the function of the form :

\[
\frac{a_0}{2} + (a_1 \, cos(x) + b_1 \, sin(x)) + ... + (a_n \, cos(x) + b_n sin(x))
\]

The coefficient ($a_0, ..., a_n, b_n$) which leads to the smallest $||f - S_n||^2$ are : 

\[
\begin{cases}
    a_n = <f, cos(nx)> = \frac{1}{\pi}\int_{-\pi}^{\pi}{f(x)\,cos(x)\,dx} \\
    b_n = <f, sin(nx)> = \frac{1}{\pi}\int_{-\pi}^{\pi}{f(x)\,sin(nx)\,dx}
\end{cases}
\] \newline

Where $S_n$ is the approximation of $f$ and Error = $f - S_n$. \newline

We get the smallest Error function with :

\[
\forall n \in \mathbb{N}, \begin{cases}
a_n = \frac{1}{\pi}\int_{-\pi}^{\pi}{f(x)\,cos(nx)\,dx} \\
b_n = \frac{1}{\pi}\int_{-\pi}^{\pi}{f(x)\,cos(nx)\,dx}
\end{cases}
\]
\subsection{Theorms}

\subsubsection{Perseval Theorms}

If $f$ is piecewise of class $\mathcal{C}^1$, then $(S_n) \overset{L_2}{\rightarrow} f$ that is, $||S_n - f|| \underset{n \rightarrow +\infty}{\rightarrow} 0$. \newline

Where : 

\[
||S_n - f|| = \sqrt{\frac{1}{\pi}\int_{-\pi}^{\pi}{(S_n(x) - f(x))^2 dx}}
\]

Consequence :

\[
f = S_n + \text{ Error} \implies ||f||^2 = ||S_n||^2 + ||\text{Error}||^2
\] \newline

Here, $||\text{Error}||^2 \underset{n \rightarrow +\infty}{\rightarrow} 0$ and $||f||^2 = \frac{1}{\pi}\int_{-\pi}^{\pi}{f^2(x)\,dx}$. \newline

Also, $||S_n||^2 = \frac{a_0^2}{2} + (a_1^2 + b_1^2) + ... + (a_n^2 + b_n^2)$ \newline

Thus, 

\[
\frac{1}{\pi}\int_{-\pi}^{\pi}{f^2(x)dx} = \frac{a_0^2}{2} + \sum_{n = 1}^{+\infty}{(a_n^2 + b_n^2)}
\]

\subsection{Dirichlet theorm}

\[
FS(f) = (S_n) \overset{\text{P.W}}{\rightarrow} \tilde{f}   
\]

Where $\tilde{f}$ is the regularized f.

\[
\forall x \in \mathbb{R}, \tilde{f}(x) = \begin{cases}
    x \text{ if f continuous at x} \\
    \frac{f(x^-) + f(x^+)}{2} \text{ if f not continuous at x}
\end{cases}
\]

\[
f = S_n + \text{ Error}
\]

\section{Complex Notation}\bigskip

\[
c_n = \frac{a_n - ib_n}{2} = \frac{1}{2\pi}\int_{-\pi}^{\pi}{f(x)e^{-inx}dx}
\]
Check :
\[
\frac{1}{2\pi}\int_{-\pi}^{\pi}{f(x)(cos(nx)-isin(nx)) \,dx} 
\]
\[
= \frac{1}{\pi}\int{f(x)\frac{cos(nx)}{2} - f(x)\frac{sin(nx)}{2}dx} = \frac{a_n}{2} - \frac{b_n}{2}
\]

Thus : 

\[
FS(f) = \sum_{n\,=\,-\infty}^{+\infty}{c_n\,e^{inx}}
\]

\newpage

We also have : 

\[
c_0\,e^{-i\times0\times x} = c_0 = \frac{a_0}{2} + \frac{b_0}{2}
\]

\[
\forall n > 0, c_{-n}e^{-inx} + c_ne^{inx} 
\]
\[
= (\frac{a_n}{2} + \frac{ib_n}{2})(cos(-nx) + isin(-nx)) + (\frac{a_n}{2} - \frac{ib_n}{2})(cos(nx) + isin(nx))
\]

Note that :
\[
C_{-n} = \overline{C_n} = \frac{a_n}{2} - i\frac{b_n}{2}
\]

\[
= \frac{a_n}{2}cos(nx) - i\frac{a_n}{2}sin(nx) + i\frac{b_n}{2}cos(nx) + \frac{b_n}{2}sin(nx) 
\]
\[
\;\;\;\;\;\;\;\; + \frac{a_n}{2}cos(nx) + i\frac{a_n}{2}sin(nx) - i\frac{b_n}{2}cos(nx) + \frac{b_n}{2}sin(nx)
\]
\[
= a_ncos(nx) + b_nsin(nx)
\]
\end{document}
